\begin{table*}[h]
\centering
\caption{Details of Catagories-Prepared Dataset}
\captionsetup[table]{skip=10pt}
\label{table:p-table}

 \def\arraystretch{0.95}
% \resizebox{\textwidth}{!}
{
\begin{tabular}{lll}
\hline
SL.No                                      & Category Name                                 & Number of Images \\ \hline
1        & Apple  & 418 \\
2        & Avocado  & 47 \\
3        & Banana  & 72 \\
4        & Kiwi  & 48 \\
5        & Lemon  & 57 \\
6        & Lime  & 42 \\
7        & Mango  & 42 \\
8        & Melon  & 163 \\
9        & Nectarine  & 51 \\
10        & Orange  & 81 \\
11        & Papaya  & 32 \\
12        & Passion-Fruit  & 28 \\
13        & Peach  & 44 \\
14        & Pear  & 140 \\
15        & Pineapple  & 29 \\
16        & Plum  & 32 \\
17        & Pomegranate  & 37 \\
18        & Red-grapefruit  & 50 \\
19        & Satsumas  & 53 \\
20        & Beverage  & 276 \\
21        & Juice  & 258 \\
22        & Milk  & 176 \\
23        & Oatghurt  & 72 \\
24        & Oat-Milk  & 51 \\
25       & Sour-Cream  & 62 \\
26        & Sour-Milk  & 36 \\
27       & Soyghurt  & 48 \\
28        & Soy-Milk  & 56 \\
29       & Yoghurt  & 85 \\
30       & Asparagus  & 14 \\
31       & Aubergine  & 28 \\
32       & Brown-Cap-Mushroom  & 25 \\
33       & Cabbage  & 19 \\
34       & Carrots  & 41 \\
35       & Cucumber  & 28 \\
36       & Garlic  & 24 \\
37       & Ginger  & 19 \\
38        & Leek  & 23 \\
39      & Onion  & 34 \\
40       & Pepper  & 103 \\
41       & Potato  & 73 \\
42       & Red-Beet  & 18 \\
43       & Tomato  & 108 \\
44       & Zucchini  & 32 \\
\hline
\end{tabular}
}
\label{table: p-result}
\end{table*}